% FILE: appendix/OC_1_3_3_K0_RESOLUTION_DISTINGUISHABILITY_PROOF.tex
% OC Core 1.3.3 scientific closure artifact.

\section{OC 1.3.3 K0 Resolution-Relative Distinguishability}
\label{app:oc133-k0-resolution}

\paragraph{Repair target.}
Core 1.3.2 used a global lower bound on raw state difference.  Core 1.3.3
replaces it with a resolution-indexed quotient.  A resolution regime is a pair
\(\rho=(\sim_\rho,\varepsilon_\rho)\), where \(\sim_\rho\subseteq S\times S\)
is an equivalence relation produced either directly by the observer/model or by
taking the equivalence closure of a weaker tolerance relation.  The positive
\(\varepsilon_\rho\) is a separation bound only on resolved equivalence
classes.

\paragraph{Definition.}
Write \(s\sim_\rho t\) when the observer or model regime cannot distinguish
\(s\) and \(t\).  The resolved substrate is \(S_\rho=S/{\sim_\rho}\), and
\([s]_\rho\neq[t]_\rho\) is permitted only when
\(\Delta_\rho([s]_\rho,[t]_\rho)\ge\varepsilon_\rho\).
The quotient difference \(\Delta_\rho\) is well-defined on classes: if
\([s]_\rho=[s']_\rho\) and \([t]_\rho=[t']_\rho\), then
\(\Delta_\rho([s]_\rho,[t]_\rho)=\Delta_\rho([s']_\rho,[t']_\rho)\).  A raw
non-transitive tolerance relation may be used only after its equivalence
closure or an explicit partition has been declared.

\paragraph{Theorem.}
Resolution-relative \(K_0\) supports continuous raw spaces without imposing
uniform discreteness on \(S\).

\paragraph{Proof.}
Let \(S=[0,1]\) with the usual metric.  No positive \(\varepsilon\) separates
all distinct raw points.  Fix \(\rho\) by partitioning \(S\) into finitely many
or locally finite equivalence cells and define \(\Delta_\rho\) on the quotient
cells by a class-level rule, for example distance between cells or an explicit
lookup table.  Because the arguments of \(\Delta_\rho\) are classes rather than
representatives, changing representatives inside the same cell does not change
the value.  Distinct quotient cells have positive separation by construction,
while points inside a cell remain raw-continuous and unresolved at \(\rho\).
Thus the threshold applies to resolved classes, not raw states.  Higher-level
continua may use the raw topology while \(K_0\) supplies only the resolved
distinguishability predicate needed for model bookkeeping.

\paragraph{Reviewer note.}
The old reading was too strong because it confused raw-state ontology with
regime-relative distinguishability.  No downstream theorem may use global
uniform discreteness unless it explicitly states a finite or discrete
realization assumption.
